\section{Controlled Checks of Placement}\label{sec:controlled-placement}
The first two deployment decisions predict different outcomes when the protected receiver or the cover executor changes. We test these predictions in a small system with inspectable boundaries. A frozen protocol precedes 72 new \texttt{gpt-6-luna} responses: 48 explanations and 24 tool plans. The same saved response or plan is replayed in each intervention arm. This pairing isolates placement from a change in model behavior. All 72 responses completed; original failed retail and negative workflow-comparison results remain in Supplements S3/S7.

\subsection{Upstream padding and a new terminal emitter}
For each of 24 public topics, a private instruction requests a brief or expanded explanation. Twelve topic pairs train the attacker and twelve disjoint pairs test it. A local origin sends each saved explanation through TLS to a frontend, which decrypts, removes padding and emits new 128-byte chunks over a separate TLS connection. Transparent TCP relays record actual ciphertext at both links. The attacker receives only the total server-to-client application-record payload length, with a threshold fitted separately for each link and arm on the training topics.

We compare no padding, upstream-only padding and padding at both links. A padded message uses sixteen 1,024-byte application frames with its length header inside encryption. TLS record parsing is independent of TCP segmentation. Every one of the 144 executions preserves the exact explanation hash and closes its socket workers.

Table~\ref{tab:controlled-placement} shows the predicted boundary change. With upstream-only padding, all upstream size observations are identical and classification falls to the balanced baseline, while the new terminal emitter retains perfect classification on the 24 held-out outputs. Padding the terminal link makes its size feature constant as well. In this pipeline, upstream equalization is lost at the new emitter and restored by padding its final link. Fixed padding costs 16,848 application-record wire bytes per link, versus a mean of 721 without padding, an aggregate byte ratio of approximately 23.4$\times$ on this workload. This counts TLS application-record headers and authentication overhead, excluding handshakes and TCP/IP headers.

This is a controlled counterexample using an intentionally length-sensitive private instruction. It establishes the re-emission mechanism and a size-only repair on this stack. Model outputs are replayed after generation; generation timing, constant-rate protection and deployed-service topic inference are separate questions. The twelve test topics provide the grouping unit; degenerate bootstrap intervals at perfect or constant prediction are retained in the materials without treating them as certainty about other workloads.

\begin{table}[t]
\centering
\caption{Controlled placement results. Panel A accuracy uses 24 held-out outputs from 12 topic pairs; byte means use all 48 outputs. Attacks are fitted separately per arm/link. Panel B replays the same 24 Luna plans. State equality is relative to each task's no-cover final state.}
\label{tab:controlled-placement}
\footnotesize
\setlength{\tabcolsep}{4pt}
\begin{tabularx}{\columnwidth}{Yrrr}
\toprule
\multicolumn{4}{l}{\textbf{A. Where padding ends}}\\
Padding & Upstream & Terminal & Terminal bytes\\
 & accuracy & accuracy & mean\\
\midrule
None & 100\% & 100\% & 721\\
Upstream only & 50\% & 100\% & 721\\
Both links & 50\% & 50\% & 16,848\\
\midrule
\multicolumn{4}{l}{\textbf{B. Where cover executes}}\\
Cover route & Task & Protected & Same\\
 & successes & violations & state\\
\midrule
No cover & 24/24 & 0/24 & 24/24\\
Discarded & 24/24 & 0/24 & 24/24\\
Isolated shadow & 24/24 & 0/24 & 24/24\\
Pure read & 24/24 & 0/24 & 24/24\\
Live write & 9/24 & 6/24 & 0/24\\
\bottomrule
\end{tabularx}
\end{table}


\subsection{Four cover routes and objective task state}
Twenty-four synthetic tasks require an inventory reservation or calendar booking, with exact references, authorized actors, explicit confirmation and protected bystander resources. Luna supplies one plan per task. The local broker validates arguments, authorization and confirmation without repairing plans. A single cover action is inserted before the original plan in four routes: discarded before dispatch, executed in a disposable state clone, executed as a pure read, or executed as a write to live state. No-cover execution supplies the paired baseline. Resource overlap, availability and protected versus spare disjoint resources are fixed before collection.

Baseline, discarded, shadow and read-only routes each complete 24/24 tasks and produce exactly the same final business state. Shadow execution performs 24 additional calls in clones; read-only cover performs 24 additional live reads. Live writes change the final state in all 24 tasks. They retain the requested goal in 15/24, but six of these violate protected-state constraints, leaving 9/24 complete successes. Relative to baseline, the paired success difference is $-62.5$ percentage points (task-bootstrap 95\% interval $[-79.2,-41.7]$). Writes can also be harmless to the stated objective: six spare-resource tasks and three slack-inventory tasks still succeed.

The distinction between goal completion and constraint preservation changes the judgment: counting only the requested booking or reservation would report 62.5\% rather than 37.5\% utility. Discard, isolation and pure reads provide concrete positive state-preservation conditions; passing the tool's actor, resource and confirmation-flag checks still leaves the goal-level check necessary. The injected cover supplies this synthetic flag, not evidence of a human separately consenting to a decoy write. Audit logs are kept separately from business state. The read-only route uses a pure getter, so its result is not a claim that every idempotent write is initially harmless. These finite plan replays measure task effects and tool work; they neither measure the cover's privacy benefit nor replace interactive-agent evaluation. Supplement S9 provides the protocol, condition-level outcomes and raw-data verification.
